% Requires graphicx. Each core figure has its own a,b panels.
% Depth figure: after the paragraph reporting fourth-block and final-output errors.
\begin{figure}[t]
\centering
\includegraphics[width=\textwidth]{figures/Fig_transformer_depth.pdf}
\caption{\textbf{Propagation of approximation error through a four-block Transformer.}
\textbf{a}, Relative error of the query-token representation after each encoder
block at $L=256$. \textbf{b}, Relative error of the query-token residual update,
defined as the difference between the representations entering and leaving
each complete encoder block. Evaluating updates measures the learned
transformation separately from the copied residual signal.
Direct transfer uses the generic HI-SMURF checkpoint from Section 3.1.
Calibration uses reference activation data to update an operator of the same
state count and TT rank while leaving Transformer weights fixed.
The analytic control uses stationary-state expectations and exact surrounding
arithmetic, including exact variances, with the same calibrated nonlinear
operators. Open markers indicate three trained-model means, each obtained
from three stream realizations for finite-stream configurations. Filled markers
and error bars denote the mean and sample standard deviation across these
three model means. Each run uses the same 128 held-out sequences.}
\label{fig:transformer_depth}
\end{figure}

% Performance figure: after the stream-length paragraph.
\begin{figure}[t]
\centering
\includegraphics[width=\textwidth]{figures/Fig_transformer_performance.pdf}
\caption{\textbf{End-to-end retrieval accuracy and numerical fidelity.}
\textbf{a}, Four-class associative-retrieval accuracy of the complete stochastic
Transformer versus bitstream length. The floating-point reference achieves
100\% accuracy on the same test sequences; chance performance is 25\%.
\textbf{b}, Relative error of the centred classification logits, evaluated
over all 128 test sequences. Subtracting each logit vector's mean removes a
common offset that does not affect class probabilities. The dashed analytic
reference is the mean over the three calibrated models.
Open markers show results averaged over three stream realizations for each
of three independently trained models. Filled markers and error bars denote
the mean and sample standard deviation across model means; stream realizations
are nested repetitions. Zero observed accuracy variation does not imply a
population confidence bound. Each stochastic phase uses $L$ sampling cycles;
the softmax and scalar FSM burn-in periods of 256 and 64 cycles, respectively,
are excluded from the horizontal axis.}
\label{fig:transformer_performance}
\end{figure}

% Ablation figure: after the final Results paragraph.
\begin{figure}[t]
\centering
\includegraphics[width=\textwidth]{figures/Fig_transformer_ablation.pdf}
\caption{\textbf{Operator coverage and simple calibration in the complete network.}
\textbf{a}, Relative centred-logit error at $L=256$ when stochastic computation
covers the attention cores, additionally all learned linear stages, and
additionally layer normalization and GELU. The final configuration evaluates
the full computational graph. All three configurations use calibrated
HI-SMURF and Sobol arithmetic; operations outside each indicated subset are
evaluated exactly. \textbf{b}, Full-graph error for direct transfer with Sobol
arithmetic, calibrated operators with independent streams, and calibrated
operators with Sobol streams. FSM driving streams remain independent
Bernoulli streams in every configuration. Faint lines connect the same trained
model across conditions. Open markers show three model means, each averaging
three stream realizations; filled markers and error bars show the mean and
sample standard deviation across these model means. No Transformer weights
are updated during calibration.}
\label{fig:transformer_ablation}
\end{figure}
